\section{\\Сравнение скоростных характеристик протоколов Bluetooth (BR, EDR, BLE)}

\subsection{Цель работы}
Провести экспериментальное сравнение реальной пропускной способности каналов связи, организованных с использованием технологий классического Bluetooth (Basic Rate и Enhanced Data Rate) и Bluetooth Low Energy. Изучить влияние параметров модуляции, размера пакетов и интервалов соединения на итоговую скорость передачи данных.

\subsection{Задачи}
\begin{itemize}[nosep]
    \item Изучить теоретические основы физических уровней Bluetooth BR, EDR и BLE: схемы модуляции, скорость символа, структуру пакетов.
    \item Рассмотреть влияние параметров L2CAP MTU, интервалов соединения и режима подтверждения (ACK/No-ACK) на пропускную способность.
    \item Подготовить стенд из двух модулей ESP32 для проведения измерений: один в роли передатчика (TX), другой — приёмника (RX).
    \item Провести серию измерений пропускной способности для трёх режимов:
    \begin{enumerate}[label=\alph*), nosep]
        \item Classic Bluetooth (BR/EDR) с использованием прошивок \texttt{ESP32\_BT\_TX.ino} и \texttt{ESP32\_BT\_RX.ino};
        \item Bluetooth Low Energy (BLE) с использованием прошивок \texttt{ESP32\_BLE\_TX.ino} и \texttt{ESP32\_BLE\_RX.ino}.
    \end{enumerate}
    \item Рассчитать среднюю скорость передачи в кбит/с и КБайт/с, построить сравнительные графики.
    \item Проанализировать полученные результаты и сформулировать выводы о применимости каждой технологии для различных сценариев передачи данных.
\end{itemize}

\subsection{Теоретический материал}

\subsubsection{Физический уровень и модуляция}
Скорость передачи данных в Bluetooth определяется схемой модуляции и символьной скоростью на физическом уровне. Во всех режимах (BR, EDR и BLE) базовая символьная скорость составляет 1~Мсимв/с, однако методы кодирования и итоговая пропускная способность существенно различаются.

Режим \textbf{Basic Rate (BR)}, соответствующий спецификации 1.x, использует гауссовскую частотную модуляцию (GFSK) с индексом модуляции $0.5$. Это обеспечивает высокую помехоустойчивость при физической скорости передачи данных 1~Мбит/с. Типичная чувствительность современных приемников в этом режиме составляет около $-90\dots-95$~дБм, что позволяет достигать стабильной связи на расстоянии до 10--30~м в помещениях.

В спецификации \textbf{Enhanced Data Rate (EDR, версия 2.0 и выше)} для увеличения пропускной способности введены дополнительные схемы фазовой модуляции. В то время как служебный заголовок пакета всегда передается с использованием GFSK, поле полезной нагрузки кодируется методами $\pi$/4-DQPSK (2~бита на символ, до 2~Мбит/с) или 8DPSK (3~бита на символ, до 3~Мбит/с). Несмотря на рост скорости, EDR сохраняет сопоставимую с BR дальность работы, хотя и предъявляет более высокие требования к отношению сигнал/шум.

Протокол \textbf{Bluetooth Low Energy (BLE, версии 4.0 и выше)} оптимизирован для сверхнизкого энергопотребления. В версиях 4.0--4.2 используется режим 1M PHY (GFSK, 1~Мбит/с). Начиная со спецификации Bluetooth 5.0, возможности физического уровня были значительно расширены:
\begin{itemize}
    \item \textbf{2M PHY:} Удваивает символьную скорость до 2~Мсимв/с, обеспечивая пиковую скорость передачи данных 2~Мбит/с.
    \item \textbf{Coded PHY (S=2 и S=8):} Вводит избыточное кодирование (FEC), что повышает чувствительность приемника до $-100\dots-105$~дБм. Это позволяет увеличить дальность связи до нескольких сотен метров (и более 1~км в условиях прямой видимости), снижая при этом полезную скорость до 500~кбит/с или 125~кбит/с соответственно.
\end{itemize}
Для повышения стабильности соединения в BLE часто применяется «стабильный индекс модуляции», варьирующийся в пределах $0.45\dots0.55$, что делает технологию оптимальной для критически важных узлов IoT.

\subsubsection{Влияние протоколов верхних уровней на пропускную способность}
Реальная скорость передачи полезных данных (throughput) всегда ниже сырой скорости физического уровня из-за служебных накладных расходов:

\begin{enumerate}[nosep]
    \item \textbf{Заголовки стека}: каждый пакет инкапсулируется в заголовки L2CAP, RFCOMM (для Classic) или ATT (для BLE), HCI, Baseband.
    \item \textbf{Подтверждения (ACK)}: в надёжных режимах (RFCOMM, BLE Indicate) передатчик ожидает ACK от приёмника, что вносит задержку и снижает эффективную скорость.
    \item \textbf{Интервал соединения (Connection Interval)}: в BLE устройства обмениваются пакетами только в определённые временные слоты. Типичные значения: 7.5--4000 мс. Чем больше интервал, тем ниже средняя скорость, но выше энергоэффективность.
    \item \textbf{MTU (Maximum Transmission Unit)}: определяет максимальный размер пакета протокола ATT. 
Для \textbf{Classic Bluetooth} через RFCOMM типичный размер кадра составляет от 255 до 1011 байт (при базовом MTU L2CAP в 672 байта), вплоть до 65535 байт.
Для \textbf{BLE} стандартный MTU по умолчанию составляет 23~байта (из которых 3~байта — заголовок ATT, оставляя 20~байт для данных). 
В современных реализациях (Bluetooth 4.2+) через процедуру \textit{MTU Exchange} это значение обычно увеличивают до 247~байт для оптимального соответствия максимальному размеру PDU канального уровня (251~байт), однако ОС Android поддерживает MTU до 517~байт.

\end{enumerate}

\subsubsection{Особенности измерений пропускной способности}
Для получения воспроизводимых результатов необходимо:
\begin{itemize}[nosep]
    \item Использовать фиксированный размер полезной нагрузки.
    \item Проводить измерения в течение достаточного интервала (не менее 10 секунд) для усреднения флуктуаций.
    \item Фиксировать параметры соединения: интервал, MTU, режим подтверждения.
\end{itemize}

\subsection{Порядок выполнения лабораторной работы}

\textbf{Оборудование:} два модуля ESP32, ПК с установленной Arduino IDE и драйверами, два окна терминала для мониторинга последовательного порта.

\textbf{Программное обеспечение:}
\begin{itemize}[nosep]
    \item Для Classic Bluetooth: прошивки \texttt{ESP32\_BT\_TX.ino} (передатчик) и \texttt{ESP32\_BT\_RX.ino} (приёмник).
    \item Для Bluetooth Low Energy: прошивки \texttt{ESP32\_BLE\_TX.ino} (передатчик) и \texttt{ESP32\_BLE\_RX.ino} (приёмник).
\end{itemize}

\begin{enumerate}[nosep, leftmargin=*]
    \item \textbf{Измерение скорости для Classic Bluetooth (BR/EDR):}
    \begin{enumerate}[label=\alph*), nosep]
        \item Загрузите прошивку \texttt{ESP32\_BT\_TX.ino} и \texttt{ESP32\_BT\_RX.ino}.
        \item В скетче убедитесь, что макрос \texttt{address[]} содержит MAC-адрес модуля с прошивкой \texttt{ESP32\_BT\_RX.ino}.
        \item Откройте два окна \texttt{Serial Monitor} (115200 бод): для TX и RX.
        \item Перезагрузите оба устройства. Дождитесь сообщения \texttt{[TX] Connected! Starting test (10 sec)...}.
        \item По окончании теста зафиксируйте вывод в окне передатчика вида:
        \begin{lstlisting}
[TX] 2048000 bytes | 10000 ms | 1638.4 kbps | 204.8 KB/s
        \end{lstlisting}
        \item Повторите измерение 3 раза, вычислите среднее значение.
    \end{enumerate}

    \item \textbf{Измерение скорости для Bluetooth Low Energy:}
    \begin{enumerate}[label=\alph*), nosep]
        \item Загрузите прошивку \texttt{ESP32\_BLE\_TX.ino} в модуль-передатчик.
        \item Убедитесь, что \texttt{serviceUUID} и \texttt{charUUID} совпадают в прошивках \texttt{ESP32\_BLE\_TX.ino} и \texttt{ESP32\_BLE\_RX.ino}.
        \item Откройте \texttt{Serial Monitor} для обоих устройств.
        \item Перезагрузите устройства. Дождитесь завершения 10-секундного теста.
        \item Зафиксируйте вывод вида:
        \begin{lstlisting}
[RX] Test complete. Received: 122880 bytes. Duration: 10.02 sec. Speed: 12.23 KB/s
        \end{lstlisting}
        \item Повторите измерение 3 раза, вычислите среднее значение.
    \end{enumerate}

    \item \textbf{Сравнительный анализ:}
    \begin{itemize}[nosep]
        \item Постройте таблицу результатов (Табл.~\ref{tab:results}) с указанием минимального, максимального и среднего значения скорости для каждого режима.
        \item Рассчитайте отношение пропускной способности EDR к BR и BLE к BR в процентах.
        \item Проанализируйте, какие факторы (размер пакета, режим подтверждения, интервал соединения) оказали наибольшее влияние на результат.
    \end{itemize}

    \item \textbf{Опционально: влияние размера пакета на скорость}
    \begin{itemize}[nosep]
        \item В прошивке \texttt{ESP32\_BLE\_TX.ino} измените размер массива \texttt{payload} с 490 байт на 20, 100 и 247 байт.
        \item Проведите измерения для каждого размера, постройте график зависимости скорости от размера пакета.
        \item Объясните наблюдаемую зависимость, опираясь на теорию накладных расходов протокола ATT.
    \end{itemize}
\end{enumerate}

\begin{table}[htbp]
\centering
\caption{Результаты измерений пропускной способности}
\label{tab:results}
\begin{tabularx}{\textwidth}{lXXXX}
\toprule
\textbf{Режим} & \textbf{Мин. скорость} & \textbf{Макс. скорость} & \textbf{Средняя скорость} & \textbf{Отклонение} \\
& \textbf{(кбит/с)} & \textbf{(кбит/с)} & \textbf{(кбит/с)} & \textbf{(\%)} \\
\midrule
Classic BR/EDR & & & & \\
BLE (1M PHY, MTU=23) & & & & \\
BLE (1M PHY, MTU=247)* & & & & \\
\bottomrule
\multicolumn{5}{l}{\small * — при поддержке согласования расширенного MTU} \\
\end{tabularx}
\end{table}

\subsection{Содержание отчета}
\begin{enumerate}[nosep, leftmargin=*]
    \item Заголовок согласно приложению.
    \item Цель работы.
    \item Задание согласно варианту.
    \item Краткое описание использованных прошивок (ссылки на файлы \texttt{ESP32\_BT\_TX.ino}, \texttt{ESP32\_BT\_RX.ino}, \texttt{ESP32\_BLE\_TX.ino}, \texttt{ESP32\_BLE\_RX.ino}).
    \item Таблица результатов измерений (Табл.~\ref{tab:results}).
    \item Графики сравнения пропускной способности (столбчатая диаграмма или линейный график).
    \item Анализ влияния параметров (размер пакета, режим подтверждения) на скорость.
    \item Ответы на контрольные вопросы.
    \item Выводы о применимости BR/EDR/BLE для различных сценариев (передача файлов, телеметрия, аудио).
\end{enumerate}

\textbf{Примечание.} При оформлении отчёта допускается приводить только ключевые фрагменты вывода \texttt{Serial Monitor}. Полные логи можно приложить отдельным файлом.

\section*{Контрольные вопросы}
\begin{enumerate}[nosep, leftmargin=*]
    \item Почему реальная пропускная способность канала всегда ниже сырой скорости физического уровня?
    \item Какие накладные расходы вносит протокол RFCOMM при передаче данных в классическом Bluetooth?
    \item Почему даже при одинаковой сырой скорости 1 Мбит/с пропускная способность BLE значительно ниже, чем у BR?
    \item Как размер пакета (payload) влияет на эффективность передачи в протоколе ATT (BLE)?
    \item В чём преимущество режима \texttt{Write Without Response} в BLE и когда его целесообразно использовать?
    \item Как интервал соединения (Connection Interval) влияет на среднюю скорость и энергопотребление в BLE?
    \item Почему в сценариях передачи больших объёмов данных (файлы, изображения) предпочтительнее использовать EDR, а не BLE?
    \item Какие факторы окружающей среды могут существенно снизить скорость передачи в беспроводных каналах?
\end{enumerate}
